\chapter{Evaluation}
\label{chapter:evaluation}

% TODO: This chapter is a methodology skeleton. Fill in testbed specifics,
% figures, and measured numbers once the experiments have been run.

This chapter evaluates the three claims made in Chapter~\ref{chapter:introduction}: that telemetry-driven feasibility filtering reduces SLA violations, that semantic contextual-bandit routing improves the accuracy-energy trade-off over resource-agnostic baselines, and that the trigger-driven allocation controller improves latency behaviour under shifting load. Each claim maps to one ablation: the full system is compared against variants with the corresponding component disabled.

\section{Experimental Setup}

\textbf{Testbed.}
% TODO: exact hardware. Placeholder description below.
The system is deployed on an edge-cloud testbed consisting of resource-constrained edge nodes and a high-capacity cloud node, connected through the shared Redis queue described in Chapter~\ref{chapter:model}. Edge nodes host the smaller models of the pool (0.5B--7B parameters), the cloud node the larger ones (12B--27B). The scheduler runs as a gateway on the edge tier.

\textbf{Workload.}
Queries are drawn from the five datasets underlying the task-type feature (MMLU, GSM8K, HellaSwag, WinoGrande, CNN/DailyMail), interleaved to create a mixed stream; ground-truth references enable online reward computation. % TODO: sample sizes, arrival process (open vs. closed loop), number of runs and seeds.

\textbf{Metrics.}
(i) SLA violation rate: fraction of queries whose end-to-end latency exceeds $T_{max}$; (ii) answer accuracy per task type; (iii) energy per query, measured at the GPU; (iv) the scalarised reward $r_t$; (v) for the allocation controller: trigger frequency, LP re-solve count, and time-to-recover after a load shift.

\textbf{Baselines.}
(a) \emph{Largest-model}: every query to the most capable cloud model; (b) \emph{smallest-model}: every query to the cheapest edge model; (c) \emph{random routing}; (d) \emph{GreenServ-style}: the contextual bandit without feasibility filtering or allocation control, i.e.\ the single-cloud assumption of prior routing work.

\section{Effect of the Feasibility Filter}
\label{sec:eval-filter}
% Ablation: full system vs. bandit-only routing (no Stage-1 filtering).
% Load is increased until edge nodes saturate; report SLA violation rate and
% tail latency for both variants.
% TODO: run and report.

\section{Accuracy-Energy Trade-off of Semantic Routing}
\label{sec:eval-routing}
% Full system vs. baselines (a)-(d) on the mixed workload; report accuracy,
% energy per query, and reward over time (learning curve).
% Sensitivity: sweep the reward weights alpha/beta and plot the resulting
% accuracy-energy Pareto front.
% TODO: run and report.

\section{Contribution of the Allocation Controller}
\label{sec:eval-allocation}
% Ablation: full system vs. system with a static model-to-node mapping.
% Scenario: abrupt load shift (e.g. workload mix flips from easy to hard
% prompts, moving demand from edge to cloud models); report latency over time,
% trigger firings, and recovery behaviour. Verify the cooldown suppresses
% oscillation.
% TODO: run and report.
